\section{Absolute residuals: controlled paired comparisons}
\label{sec:absolute}

\paragraph{Data generation and fitted prediction rule.}
Each trial independently draws 800 training observations, $n$ calibration
observations, and 1,200 test observations. Let $X\sim N(0,I_3)$,
$g(X)=\mathbf 1\{X_1>0\}$, and, for $j=1,\ldots,d$,
\begin{equation}\label{eq:toy-model}
 Y_j=X^\top\beta_j+(1+0.8g(X))\tau_j\varepsilon_j,
 \quad \beta_{kj}=\tfrac12\sin\{(k-1)d+j\},\quad
 \tau_j=0.6\,4^{(j-1)/(d-1)}.
\end{equation}
The four noise laws have zero mean and unit variance: independent Gaussian;
Gaussian with equicorrelation $0.8$; independent Laplace with scale
$1/\sqrt2$; and independent standardized lognormal,
$\varepsilon_j=(e^{0.8Z_j}-e^{0.32})/
\sqrt{(e^{0.64}-1)e^{0.64}}$ with $Z_j\sim N(0,1)$.
An OLS model with intercept is fitted on training observations only. Both
methods use the same absolute residuals
$E_j(x,y_j)=|y_j-\widehat f_j(x)|$, calibration data, and test observations.
Thus any difference in this section comes from calibration, not refitting
or changing the residual definition. A threshold vector $L$ gives
$\prod_j[\widehat f_j(x)-L_j,\widehat f_j(x)+L_j]$ with volume
$2^d\prod_jL_j$.

\paragraph{Scope and uncertainty.}
The calibration-size study uses $d=6$, $\alpha=0.1$, $n\in\{30,80,200\}$,
four noise laws, and 120 independent fitted trials per setting (1,440 total).
The target study keeps $d=6,n=80$ and uses Gaussian and Laplace noise at
$\alpha\in\{0.05,0.1,0.2,0.3\}$; the $0.1$ results are the same trials as
in the first study. There are 2,160 distinct absolute-score trials overall.
These are retained independent fitted trials from the earlier report, with
both methods rerun from the saved scores for this revision; they are not
new independent replications of that evidence.

Joint coverage is the fraction of test observations inside all coordinate
intervals. Size comparisons average the \emph{within-trial} volume ratio
$V_{\rm env}/V_{\rm old}$, not a ratio of pooled means. Error bars are
pointwise 95\% Student-$t$ Monte Carlo intervals across the 120 trial values.
Test observations within one fitted trial are not treated as independent
replications for these error bars. All absolute regions here are finite.
The old implementation retains its original deterministic $10^{-10}$ score
jitter; the mathematical theorem uses closed cells without jitter. The
reported numerical checks use a relative endpoint tolerance of $10^{-8}$.

\begin{figure}[tbp]
\centering\includegraphics[width=\linewidth]{absolute_n.pdf}
\caption{Calibration-size comparison on identical absolute residuals.
Top: joint coverage, with target $0.9$ dashed. Bottom: paired volume ratios;
faint points are individual fitted trials, diamonds are means, and bars are
95\% Monte Carlo intervals. Each setting has 120 trials, $d=6$.}
\label{fig:absolute-n}
\end{figure}

\paragraph{Calibration size: the strongest gains occur with limited data.}
At $n=30$, the mean paired volume reductions are 5.64\% (Gaussian),
3.11\% (correlated Gaussian), 7.36\% (Laplace), and 16.44\% (skewed
lognormal). Their 95\% Monte Carlo half-widths are respectively 0.67,
0.34, 1.01, and 2.32 percentage points. At $n=80$ the reductions narrow
to 0.77--1.37\%; at $n=200$ they are 0.23--0.34\%.
The envelope therefore offers a distribution-free direction of improvement,
but its magnitude is strongly setting dependent. The skewed, small-sample
case is the clearest empirical gain in this grid.

Across all 2,160 paired trials and 12,960 coordinate comparisons, there are
no containment violations. In the size study, envelope mean joint coverage
ranges from 0.9006 to 0.9135 and stays close to the 0.9 target. It is slightly
lower than old-shortcut coverage, as nesting requires; shrinking a conservative
region is the intended efficiency gain. The experiment supports the theorem
but does not substitute for it.

\paragraph{Different target levels.}
At $n=80$, paired volume reductions span 1.05--1.64\% for Gaussian noise
and 0.85--1.85\% for Laplace across the four targets. The largest reductions
in both families occur at $\alpha=0.05$, but the curves are not generally
monotone in $\alpha$. Mean envelope coverage exceeds its nominal target by
0.03--1.15 percentage points across these settings. Thus the refinement
persists across targets without a universal claim about how its magnitude
must change with $\alpha$.


\begin{figure}[tbp]
\centering\includegraphics[width=0.88\linewidth]{absolute_alpha.pdf}
\caption{Miscoverage sensitivity at $n=80,d=6$, 120 fitted trials per point.
Top: coverage minus $1-\alpha$ (zero is nominal). Bottom: paired volume ratio.
Error bars describe Monte Carlo uncertainty, not a simultaneous coverage test.}
\label{fig:absolute-alpha}
\end{figure}

\paragraph{What the regions look like.}
Figure~\ref{fig:absolute-region} uses a separately generated two-dimensional
Laplace trial, $n=30$, $\alpha=0.1$, seed 2026990918 and test input 0, fixed
before viewing the geometry. The displayed axes subtract the fitted center
and divide by the known DGP scales $\tau_j$, solely for visualization.
The enlargement near the upper-right corner makes a modest boundary change
visible while the main axes preserve its true size. The area ratio is
$0.9750$: a 2.50\% reduction. Gray shading evaluates the full-conformal
acceptance rule on a $301\times301$ grid; every accepted grid point is
inside the envelope. This grid is an illustration, not an exact set-area
calculation or a coverage experiment. Full LWC is reviewed mathematically;
the report does not run its exponential full-cell enumeration.

\begin{figure}[H]
\centering\includegraphics[width=\linewidth]{absolute_region.pdf}
\caption{One fixed-input absolute-residual region. Orange dashed: old
shortcut; teal: envelope; gray: full-conformal grid acceptance. Both axes
use the same coordinate normalization for every method. The right panel shows the
upper-right corner without changing the scale of the main plot.}
\label{fig:absolute-region}
\end{figure}
